\lecture{Искусство периода феодальной раздробленности и начала объединения русских земель}{lecture-04}

{
  \title{История Русского Искусства}
  \subtitle{Лекция 4: объединение русских земель}
  \date{22 сентября 2026}

  \begin{frame}[plain,noframenumbering]
    \titlepage
  \end{frame}
}

\sectionimage{lectures/03-westrussian/images/section-logos/tower.png}
\section{Искусство позднего Средневековья}

\begin{frame}{После нашествия}
  К периоду монголо-татарского нашествия развитие искусства русских земель
  достигло высокого уровня: архитектура, живопись, скульптура,
  прикладное искусство, миниатюра и художественное ремесло.

  \vspace{0.8em}
  Естественное движение искусства и культуры было прервано в тот момент,
  когда перед ним раскрывались широкие перспективы.
\end{frame}

\begin{frame}{Начало объединения}
  В долгой борьбе за руководство общерусским движением между Тверью
  и Москвой роль лидера досталась московским князьям.

  \begin{itemize}
    \item 1378~г. --- первое поражение Орды на реке Воже.
    \item 1380~г. --- Куликовская битва.
    \item Национальный подъем и идея единства становятся важной
          предпосылкой нового художественного развития.
  \end{itemize}
\end{frame}

\begin{frame}{Русские земли XIV века}
  \widefigure[Wikimedia Commons]
    {lectures/04-feodalreunion/images/maps/russian-lands-map.jpg}
    {Русские земли в конце XIV~в.:
     художественные школы развиваются на фоне соперничества Москвы,
     Твери, Новгорода, Пскова и Орды}
\end{frame}

\subsection{Новгород XIV--XV вв.}

\begin{frame}{Новгород после Куликовской битвы}
  Усиление избранного новгородского войска под командованием воеводы
  Ивана Мошкова способствовало победе на Куликовом поле.

  \begin{itemize}
    \item Но концентрация русских земель вокруг Москвы постепенно
          сокращала политические возможности Новгорода.
    \item С XV~в. в новгородской культуре ярко проявляется
          консервативная ориентация на искусство прошлого.
    \item Идеализируется новгородская старина.
  \end{itemize}
\end{frame}

\begin{frame}{Городская среда}
  Новгород оставался городом острых социальных конфликтов
  и незатухавшей классовой борьбы.

  \begin{itemize}
    \item Здесь оформляются первые массовые городские еретические
          движения.
    \item Они подрывали устои церкви и догматические основы православия.
    \item Эта среда формирует сложную, противоречивую и динамичную
          картину искусства XIV--XV~вв.
  \end{itemize}
\end{frame}

\subsection{Архитектура Новгорода}

\begin{frame}{Новые принципы}
  В архитектуре Новгорода с конца XII~в. нарастают важные изменения.

  \begin{itemize}
    \item Местный волховский плитняк, валуны и кирпич формируют
          неповторимое пластическое звучание.
    \item Полуцилиндрические своды меняют верх и внешний облик храма.
    \item Одна апсида часто заменяет трехапсидную композицию.
    \item Фасадная декорация становится самостоятельнее конструкции.
  \end{itemize}
\end{frame}

\begin{frame}{Спас на Ковалеве}
  \widefigure[Wikimedia Commons]
    {lectures/04-feodalreunion/images/novgorod-architecture/kovalevo.jpg}
    {Церковь Спаса на Ковалеве, 1342~г.:
     на старой композиционной основе возникает новый городской тип храма}
\end{frame}

\begin{frame}{Успение на Волотовом поле}
  \tallfigure[Wikimedia Commons]
    {lectures/04-feodalreunion/images/novgorod-architecture/volotovo.jpg}
    {Церковь Успения на Волотовом поле, 1352~г.:
     камерность и простота становятся частью новгородского
     архитектурного языка}
\end{frame}

\begin{frame}{Феодор Стратилат на Ручью}
  \widefigure[Wikimedia Commons]
    {lectures/04-feodalreunion/images/novgorod-architecture/feodor-stratilat.jpg}
    {Церковь Феодора Стратилата на Ручью, 1360--1361~гг.:
     расцвет новгородской архитектуры XIV~в.}
\end{frame}

\begin{frame}{Законченный образ}
  В церкви Феодора Стратилата особенно важно фасадное оформление.

  \begin{itemize}
    \item Лопатки соответствуют местам внутренних столбов.
    \item Многолопастные арки подчеркивают характер перекрытия.
    \item Стрельчатые перспективные порталы и глубокие окна
          делают стену живой пластической поверхностью.
    \item Этот образ надолго становится ориентиром для школы.
  \end{itemize}
\end{frame}

\begin{frame}{Спас на Ильине улице}
  \widefigure[Wikimedia Commons]
    {lectures/04-feodalreunion/images/novgorod-architecture/spas-ilyina.jpg}
    {Церковь Спаса Преображения на Ильине улице, 1374~г.:
     декор начинает соперничать с конструкцией}
\end{frame}

\begin{frame}{Поздние варианты типа}
  \begin{columns}[c,onlytextwidth]
    \begin{column}{0.48\textwidth}
      \centering
      \includegraphics[width=\linewidth,height=0.62\textheight,keepaspectratio]
        {lectures/04-feodalreunion/images/novgorod-architecture/john-vitka.jpg}\\[0.5em]
      \figcaption{Иоанн Богослов на Витке, 1384~г.}
    \end{column}
    \begin{column}{0.48\textwidth}
      \centering
      \includegraphics[width=\linewidth,height=0.62\textheight,keepaspectratio]
        {lectures/04-feodalreunion/images/novgorod-architecture/peter-paul-kozhevniki.jpg}\\[0.5em]
      \figcaption{Петр и Павел в Кожевниках, 1406~г.}
    \end{column}
  \end{columns}
\end{frame}

\begin{frame}{Итог архитектуры}
  Поиски выразительности в пределах сложившегося типа продолжаются.
  Но архитектура уже не выдвигает принципиально новых идей.

  \begin{itemize}
    \item Она ориентируется не только на XIV~в., но и на XII~в.
    \item Это связано с усилением Москвы.
    \item Одновременно углубляется новгородский сепаратизм,
          ассоциированный с «золотым веком» города.
  \end{itemize}
\end{frame}

\begin{frame}{Евфимий II}
  Архиепископ Евфимий II Брадатый служил около тридцати лет
  и поставил Новгород в центр русского исторического процесса.

  \vspace{0.8em}
  Каменное строительство он понимал как демонстрацию силы
  и могущества города.
\end{frame}

\begin{frame}{Владычная палата}
  \widefigure[Wikimedia Commons]
    {lectures/04-feodalreunion/images/novgorod-architecture/vladychnaya-chamber.jpg}
    {Владычная / Грановитая палата в Новгородском детинце,
     1433~г.: редкий памятник гражданской архитектуры}
\end{frame}

\begin{frame}{Готические своды}
  \tallfigure[Wikimedia Commons]
    {lectures/04-feodalreunion/images/novgorod-architecture/vladychnaya-vault.jpg}
    {Свод Грановитой палаты:
     светская архитектура входит в новгородскую репрезентацию}
\end{frame}

\subsection{Феофан Грек}

\begin{frame}{Золотой век живописи}
  XIV~в. --- золотой век новгородской монументальной живописи.
  Местная школа работает в тесном контакте с приезжими мастерами
  византийско-славянского мира.

  \vspace{0.8em}
  Первый среди них --- Феофан Грек, мастер из Константинополя
  (30-е гг. XIV~в. -- после 1405~г.).
\end{frame}

\begin{frame}{Фрески Феофана}
  \tallfigure[Wikimedia Commons]
    {lectures/04-feodalreunion/images/novgorod-frescoes/spas-ilyina-pantocrator.jpg}
    {Христос Пантократор.
     Фреска церкви Спаса Преображения на Ильине улице}
\end{frame}

\begin{frame}{Столпники}
  \tallfigure[Wikimedia Commons]
    {lectures/04-feodalreunion/images/novgorod-frescoes/spas-ilyina-stylites.jpg}
    {Феофановская манера:
     стремительный рисунок, резкие высветления и духовное напряжение}
\end{frame}

\begin{frame}{Старцы Феофана Грека}
  В барабане Спаса на Ильине Феофан собирает не обычный
  пророческий ряд, а ветхозаветных праведников и пророков:
  Адама, Авеля, Сифа, Еноха, Ноя, Мельхиседека, Илию
  и Иоанна Предтечу.

  \begin{itemize}
    \item Адам --- чистота и мудрость до грехопадения;
          Авель --- активный носитель добра.
    \item Сиф связан с передачей людям грамоты и знания;
          Енох --- праведник, угодивший Богу и взятый от смерти.
    \item Ной становится новым началом человеческого рода;
          Мельхиседек --- «царь правды».
    \item Илья являет нестяжание и пророческую силу;
          жизнь Иоанна Предтечи прочитана как великий пост.
  \end{itemize}

  Это собирательный образ редкой мощи, воли и самоограничения. Таков праведный идеал Феофана.
\end{frame}

\begin{frame}{Страстная драматичность}
  В системе росписей важны страстные сюжеты.

  \begin{itemize}
    \item Бог Феофана --- грозная, почти непостижимая сила.
    \item Человек рядом с ней чувствует собственные пределы.
    \item Фигуры обладают трагической психологической глубиной.
    \item Манера письма фиксирует драму, а не спокойное повествование.
  \end{itemize}
\end{frame}

\begin{frame}{Влияние Феофана}
  Отголоски влияния Феофана заметны во фресках церкви Феодора
  Стратилата и церкви Успения на Волотовом поле.

  \begin{itemize}
    \item Есть напряженное стилистическое родство.
    \item Но палитра становится ярче и праздничнее.
    \item Композиция мягче, а драматизм менее обнажен.
  \end{itemize}
\end{frame}

\subsection{Новгородская икона}

\begin{frame}{Черты новгородской иконы}
  В XIV--XV~вв. постепенно складывается яркий и самобытный стиль.

  \begin{itemize}
    \item Простая и выразительная плоскостная композиция.
    \item Чеканный, обобщенный силуэт.
    \item Чистая, звонкая, высветленная палитра.
    \item Коренастые и крепкие фигуры.
    \item Лаконичный и точный изобразительный язык.
  \end{itemize}
\end{frame}

\begin{frame}{Георгий с житием}
  \tallfigure[Google Art Project; Wikimedia Commons]
    {lectures/04-feodalreunion/images/novgorod-icons/saint-george-life.jpg}
    {Чудо Георгия о змие с житием,
     Новгород, первая половина XIV~в.}
\end{frame}

\begin{frame}{Святой Георгий}
  \tallfigure[Wikimedia Commons]
    {lectures/04-feodalreunion/images/novgorod-icons/george-15c.jpg}
    {Святой Георгий, новгородская школа, XV~в.:
     воинский образ превращается в знак победы и защиты}
\end{frame}

\begin{frame}{Избранные святые}
  Во второй половине XIV~в. популярны изображения избранных святых:
  по двое, трое, четверо и больше; во весь рост и по пояс.

  \begin{itemize}
    \item Важны фронтальность и статуарность.
    \item Различаются атрибуты, жесты и постановка рук.
    \item Святые отвечают потребностям города:
          защита от пожаров, покровительство торговле,
          земледелию и скоту.
  \end{itemize}
\end{frame}

\begin{frame}{Битва новгородцев с суздальцами}
  \widefigure[Третьяковская галерея; Wikimedia Commons]
    {lectures/04-feodalreunion/images/novgorod-icons/battle-novgorod-suzdal-detail.jpg}
    {Фрагмент иконы «Битва новгородцев с суздальцами»:
     историческая память становится образом небесного покровительства}
\end{frame}

\begin{frame}{Флор и Лавр}
  \tallfigure[Wikimedia Commons]
    {lectures/04-feodalreunion/images/novgorod-icons/miracle-florus-laurus.jpg}
    {«Чудо о Флоре и Лавре»:
     покровительство коням, хозяйству и повседневной городской жизни}
\end{frame}

\begin{frame}{Северные письма}
  Иконы местных провинциальных мастеров условно называют
  «северными письмами».

  \begin{itemize}
    \item Упрощенная композиция.
    \item Угловатый, жестковатый рисунок.
    \item Наивная суровая выразительность.
    \item Связь христианского образа с простоватым
          мифологическим персонажем.
    \item Контраст мягких бледных и плотных цветов.
  \end{itemize}
\end{frame}

\begin{frame}{Поздняя иконопись}
  В XVI~в. живописность постепенно уступает иконности,
  суховатой правильности форм.

  \vspace{0.8em}
  Но сохраняются произведения, где заимствование у Москвы
  соединяется с местной особенностью и оригинальностью.
\end{frame}

\subsection{Псков}

\begin{frame}{Архитектура Пскова}
  Со второй половины XIII~в. Псков возводит прежде всего
  оборонительные сооружения.

  \begin{itemize}
    \item Город стоит на северо-западных границах Руси.
    \item Каменная стена, башни, суровые и лаконичные формы
          отвечают военной функции.
    \item Псковская пластика грубовата, но цельна и выразительна.
  \end{itemize}
\end{frame}

\begin{frame}{Троицкий собор}
  \widefigure[Wikimedia Commons]
    {lectures/04-feodalreunion/images/pskov/trinity-cathedral.jpg}
    {Псковский Кром и Троицкий собор:
     поздний облик главного городского сакрального центра}
\end{frame}

\begin{frame}{Псковская живопись}
  Псковская живопись XIV--XV~вв. очень оригинальна и подвижна.

  \begin{itemize}
    \item Ей свойственны напряженные иконографические решения.
    \item Важны резкие жесты и драматическая выразительность.
    \item Палитра тяготеет к красно-коричневым и желтым тонам.
    \item Образ часто конкретнее и эмоциональнее новгородского.
  \end{itemize}
\end{frame}

\begin{frame}{Снетогорский монастырь}
  \tallfigure[Wikimedia Commons]
    {lectures/04-feodalreunion/images/pskov/snetogorsky-cathedral.jpg}
    {Собор Рождества Богородицы Снетогорского монастыря:
     главный ранний памятник псковской монументальной школы}
\end{frame}

\begin{frame}{Снетогорские фрески}
  Собор Рождества Богородицы Снетогорского монастыря был расписан
  в 1313~г.

  \begin{itemize}
    \item Художественный язык архаизирующий и прямолинейный.
    \item Фигуры одновременно человеческие, движущиеся
          и схематичные.
    \item Характерны грубоватая экспрессия, свободная трактовка
          иконографии и апокрифические мотивы.
  \end{itemize}
\end{frame}

\begin{frame}{Мелетовские фрески}
  Мелетовские фрески 1465~г. соединяют псковскую образность
  с новой техникой исполнения.

  \begin{itemize}
    \item Их сюжеты и образы сложны для расшифровки.
    \item В них сохраняется драматическая энергия псковской школы.
    \item Но заметны признаки позднего этапа: схематизм,
          стандартизация и усиление внешней выразительности.
  \end{itemize}
\end{frame}

\begin{frame}{Псков к середине XV века}
  К середине XV~в. окончательно складывается псковская школа.

  \begin{itemize}
    \item Схематизм и драматизм.
    \item Повышенная эмоциональность.
    \item Сильное выражение лица.
    \item Ослабление идеи при стандартизации композиции.
  \end{itemize}
\end{frame}

\begin{frame}{Итоги}
  Новгород и Псков XIV--XV~вв. показывают, что после разрыва XIII~в.
  русское искусство не просто возобновляется, а обретает новые
  городские и региональные формы.

  \vspace{0.8em}
  Новгород дает зрелую архитектуру малого городского храма,
  мощную монументальную живопись и яркий иконописный язык.
  Псков дает пограничную, оборонительную и драматическую школу
  с особой пластикой и экспрессией.
\end{frame}

\begin{frame}{Переход к объединению}
  На фоне политического подъема Москвы местные школы не исчезают сразу.
  Они продолжают работать, спорить с прошлым и идеализировать собственную
  старину.

  \vspace{0.8em}
  Но постепенно эти школы входят в новый общерусский художественный
  процесс, связанный с объединением земель.
\end{frame}
